% common/tex/preamble.tex
% Copyright (C) 2026 Hirokazu Maesaka (RIKEN SPring-8 Center)
% SPDX-License-Identifier: CC-BY-4.0
%
% Shared preamble for all topic documents. Loaded by topics/<topic>/main.tex with \input.
% Build with lualatex; the language is selected with \memolang (en or ja) through latexmk -pretex.

\documentclass[11pt,a4paper]{article}

% Language switch. The default is English. latexmk -pretex='\def\memolang{ja}' overrides it.
\providecommand{\memolang}{en}
% \ifx compares macro meanings. Unlike \ifdefstring, it is expandable, so \en and \ja also work inside \edef.
\def\memolangEN{en}
\def\memolangJA{ja}
\newcommand{\en}[1]{\ifx\memolang\memolangEN #1\fi}
\newcommand{\ja}[1]{\ifx\memolang\memolangJA #1\fi}

\usepackage{luatexja-preset}
\usepackage{etoolbox}
\usepackage{amsmath,amssymb}
\usepackage[style=numeric,backend=biber,sorting=none]{biblatex}
\addbibresource{common/references/references.bib}
% Memo numbers (YYYY-NNN) of the topics, generated from common/documents.toml. Cite a memo with \cite{YYYY-NNN}.
\addbibresource{common/documents.bib}

% Memo number. Each topic declares its number with \setmemonumber{YYYY-NNN} (see common/documents.toml).
% The number is shown in the page header and is used in the names of the output files.
\providecommand{\memonumber}{}
\newcommand{\setmemonumber}[1]{\def\memonumber{#1}}
\usepackage{fancyhdr}
\setlength{\headheight}{14pt}
\pagestyle{fancy}
\fancyhf{}
\fancyhead[R]{Memo~\memonumber}
\fancyfoot[C]{\thepage}
\fancypagestyle{plain}{%
  \fancyhf{}%
  \fancyhead[R]{Memo~\memonumber}%
  \fancyfoot[C]{\thepage}%
  \renewcommand{\headrulewidth}{0pt}}

\setlength{\parindent}{0pt}
\setlength{\parskip}{0.6em}

% Slide markers (write build/<topic>/<jobname>.slides for common/slides/build_slides.py).
% common/tex/slides.tex
% Copyright (C) 2026 Hirokazu Maesaka (RIKEN SPring-8 Center)
% SPDX-License-Identifier: CC-BY-4.0
%
% Slide markers. Each marker writes one line "kind:payload" to build/<topic>/<jobname>.slides during compilation.
% common/slides/build_slides.py reads that file and builds the PowerPoint deck. The markers print nothing in the PDF.
% Arguments are expanded at the point of use, so \en{...} and \ja{...} select the language of the current build.
% Do not use fragile commands (for example \cite or \label) inside slide arguments.

\newwrite\slidesfile
\AtBeginDocument{\immediate\openout\slidesfile=\jobname.slides\relax}
\AtEndDocument{\immediate\closeout\slidesfile}

% #1: entry kind, #2: payload. The payload is expanded, then written as plain characters.
\newcommand{\slideentry}[2]{%
  \begingroup
  \edef\slidetmp{#2}%
  \edef\slidedet{\detokenize\expandafter{\slidetmp}}%
  \immediate\write\slidesfile{#1:\slidedet}%
  \endgroup}

% Raw entries are written without expansion. Use them for math and paths, which contain fragile commands.
\newcommand{\slideentryraw}[2]{%
  \immediate\write\slidesfile{#1:\detokenize{#2}}}

\newcommand{\slidetitle}[1]{\slideentry{title}{#1}}
\newcommand{\slidepoint}[1]{\slideentry{point}{#1}}
\newcommand{\slidemath}[1]{\slideentryraw{math}{#1}}
\newcommand{\slidefigure}[1]{\slideentryraw{figure}{#1}}
\newcommand{\slidenote}[1]{\slideentry{note}{#1}}

